\Chapter{26}

\Verse{1}
{
    \zA{话说宝玉养过了三十三天之后，不但身体强壮，亦且连脸上疮痕平复，仍回大 观园去。}
}
{
    \eA{After thirty days' careful nursing, Pao-yü, we will now notice, not only got
        strong and hale in body, but the scars even on his face completely healed
        up;}
}
{
}

\Verse{2}
{
    \zA{这也不在话下。}
}
{
    \eA{so he was able to shift his quarters again into the garden of Broad Vista.}
}
{
}

\Verse{3}
{
    \zA{且说近日宝玉病的时节，贾芸带着家下小厮坐更看守，昼夜在这里，那小红同 众丫鬟也在这里守着宝玉。}
}
{
    \eA{But we will banish this topic as it does not deserve any additional
        explanations. Let us now turn our attention elsewhere. During the time that
        Pao-yü was of late laid up in bed, Chia Yün along with the young pages of
        the household sat up on watch to keep an eye over him,}
}
{
}

\Verse{4}
{
    \zA{彼此相见日多，渐渐的混熟了。}
}
{
    \eA{and both day and night, they tarried on this side of the mansion.}
}
{
}

\Verse{5}
{
    \zA{小红见贾芸手里拿着块 绢子，倒像是自己从前掉的，待要问他，又不好问。}
}
{
    \eA{But Hsiao Hung as well as all the other waiting-maids remained in the same
        part to nurse Pao-yü, so (Chia Yün) and she saw a good deal of each other on
        several occasions,}
}
{
}

\Verse{6}
{
    \zA{不料那和尚道士来过，用不着 一切男人，贾芸仍种树去了；这件事待放下又放不下，待要问去又怕人猜疑。}
}
{
    \eA{and gradually an intimacy sprung up between them. Hsiao Hung observed that
        Chia Yün held in his hand a handkerchief very much like the one she herself
        had dropped some time ago and was bent upon asking him for it, but she did,
        on the other hand,}
}
{
}

\Verse{7}
{
    \zA{正是 犹豫不决、神魂不定之际，忽听窗外问道：“姐姐在屋里没有？”}
}
{
    \eA{not think she could do so with propriety. The unexpected visit of the bonze
        and Taoist priest rendered, however, superfluous the services of the various
        male attendants,}
}
{
}

\Verse{8}
{
    \zA{小红闻听，在窗 眼内望外一看，原来是本院的个小丫头佳蕙，因答说：“在家里呢，你进来罢。”}
}
{
    \eA{and Chia-yün had therefore to go again and oversee the men planting the
        trees. Now she had a mind to drop the whole question, but she could not
        reconcile herself to it; and now she longed to go and ask him about it,}
}
{
}

\Verse{9}
{
    \zA{佳蕙听了跑进来，就坐在床上，笑道：“我好造化!才在院子里洗东西，宝玉叫往 林姑娘那里送茶叶，花大姐姐交给我送去。}
}
{
    \eA{but fears rose in her mind lest people should entertain any suspicions as to
        the relations that existed between them. But just as she faltered, quite
        irresolute, and her heart was thoroughly unsettled, she unawares heard some
        one outside inquire:}
}
{
}

\Verse{10}
{
    \zA{可巧老太太给林姑娘送钱来，正分给他 们的丫头们呢，见我去了，林姑娘就抓了两把给我。}
}
{
    \eA{"Sister, are you in the room or not?" Hsiao Hung, upon catching this
        question, looked out through a hole in the window; and perceiving at a
        glance that it was no one else than a young servant-girl,}
}
{
}

\Verse{11}
{
    \zA{也不知是多少，你替我收着。”}
}
{
    \eA{attached to the same court as herself, Chia Hui by name, she consequently
        said by way of reply:}
}
{
}

\Verse{12}
{
    \zA{便把手绢子打开，把钱倒出来交给小红。}
}
{
    \eA{"Yes, I am; come in!" When these words reached her ear, Chia Hui ran in, and
        taking at once a seat on the bed,}
}
{
}

\Verse{13}
{
    \zA{小红就替他一五一十的数了收起。}
}
{
    \eA{she observed with a smile: "How lucky I've been! I was a little time back in
        the court washing a few things,}
}
{
}

\Verse{14}
{
    \zA{佳蕙道：“你这两日心里到底觉着怎么样？}
}
{
    \eA{when Pao-yü cried out that some tea should be sent over to Miss Lin, and
        sister Hua handed it to me to go on the errand.}
}
{
}

\Verse{15}
{
    \zA{依我说，你竟家去住两日，请一个 大夫来瞧瞧，吃两剂药，就好了。”}
}
{
    \eA{By a strange coincidence our old lady had presented some money to Miss Lin
        and she was engaged at the moment in distributing it among their
        servant-girls. As soon therefore as she saw me get there, Miss Lin forthwith
        grasped two handfuls of cash and gave them to me;}
}
{
}

\Verse{16}
{
    \zA{小红道：“那里的话？}
}
{
    \eA{how many there are I don't know, but do keep them for me!"}
}
{
}

\Verse{17}
{
    \zA{好好儿的，家去做什么？”}
}
{
    \eA{Speedily then opening her handkerchief, she emptied the cash. Hsiao Hung
        counted them for her by fives and tens at a time.}
}
{
}

\Verse{18}
{
    \zA{佳蕙道：“我想起来了。}
}
{
    \eA{She was beginning to put them away, when Chia Hui remarked: "How are you,}
}
{
}

\Verse{19}
{
    \zA{林姑娘生的弱，时常他吃药，你就和他要些来吃，也是一 样。”}
}
{
    \eA{after all, feeling of late in your mind? I'll tell you what; you should
        really go and stay at home for a couple of days. And were you to ask a
        doctor round and to have a few doses of medicine you'll get all right at
        once!"}
}
{
}

\Verse{20}
{
    \zA{小红道：“胡说，药也是混吃的？”}
}
{
    \eA{"What are you talking about?" Hsiao Hung replied. "What shall I go home for,}
}
{
}

\Verse{21}
{
    \zA{佳蕙道：“你这也不是个长法儿，又懒 吃懒喝的，终久怎么样？”}
}
{
    \eA{when there's neither rhyme nor reason for it!" "Miss Lin, I remember, is
        naturally of a weak physique, and has constantly to take medicines,"}
}
{
}

\Verse{22}
{
    \zA{小红道：“怕什么？}
}
{
    \eA{Chia Hui added, "so were you to ask her for some and bring them over and
        take them,}
}
{
}

\Verse{23}
{
    \zA{还不如早些死了倒干净。”}
}
{
    \eA{it would come to the same thing." "Nonsense!" rejoined Hsiao Hung,}
}
{
}

\Verse{24}
{
    \zA{佳蕙道： “好好儿的，怎么说这些话？”}
}
{
    \eA{"are medicines also to be recklessly taken ?" "You can't so on for ever like
        this,"}
}
{
}

\Verse{25}
{
    \zA{小红道：“你那里知道我心里的事！”}
}
{
    \eA{continued Chia Hui; "you're besides loth to eat and loth to drink, and what
        will you be like in the long run?"}
}
{
}

\Verse{26}
{
    \zA{佳蕙点头， 想了一会道：“可也怨不得你。}
}
{
    \eA{"What's there to fear?" observed Hsiao Hung; "won't it anyhow be better to
        die a little earlier? It would be a riddance!"}
}
{
}

\Verse{27}
{
    \zA{这个地方，本也难站。}
}
{
    \eA{"Why do you deliberately come out with all this talk?"}
}
{
}

\Verse{28}
{
    \zA{就像昨儿老太太因宝玉病了 这些日子，说伏侍的人都辛苦了，如今身上好了，各处还香了愿，叫把跟着的人都 按着等儿赏他们。}
}
{
    \eA{Chia Hui demurred. "How could you ever know anything of the secrets of my
        heart?" Hsiao Hung inquired. Chia Hui nodded her head and gave way to
        reflection. "I don't think it strange on your part," she said after a time;
        "for it is really difficult to abide in this place! Yesterday, for instance,}
}
{
}

\Verse{29}
{
    \zA{我们算年纪小，上不去，我也不抱怨；像你怎么也不算在里头？}
}
{
    \eA{our dowager lady remarked that the servants in attendance had had, during
        all the days that Pao-yü was ill, a good deal to put up with, and that now
        that he has recovered,}
}
{
}

\Verse{30}
{
    \zA{我心里就不服。}
}
{
    \eA{incense should be burnt everywhere,}
}
{
}

\Verse{31}
{
    \zA{袭人那怕他得十分儿，也不恼他，原该的。}
}
{
    \eA{and the vows fulfilled; and she expressed a wish that those in his service
        should, one and all, be rewarded according to their grade.}
}
{
}

\Verse{32}
{
    \zA{说句良心话，谁还能比 他呢？}
}
{
    \eA{I and several others can be safely looked upon as young in years,}
}
{
}

\Verse{33}
{
    \zA{别说他素日殷勤小心，就是不殷勤小心，也拼不得。}
}
{
    \eA{and unworthy to presume so high; so I don't feel in any way aggrieved; but
        how is it that one like you couldn't be included in the number?}
}
{
}

\Verse{34}
{
    \zA{只可气晴雯绮霞他们这 几个都算在上等里去，伏着宝玉疼他们，众人就都捧着他们。}
}
{
    \eA{My heart is much annoyed at it! Had there been any fear that Hsi Jen would
        have got ten times more, I could not even then have felt sore against her,
        for she really deserves it!}
}
{
}

\Verse{35}
{
    \zA{你说可气不可气？”}
}
{
    \eA{I'll just tell you an honest truth; who else is there like her?}
}
{
}

\Verse{36}
{
    \zA{小红道：“也犯不着气他们。}
}
{
    \eA{Not to speak of the diligence and carefulness she has displayed all along,}
}
{
}

\Verse{37}
{
    \zA{俗语说的：‘千里搭长棚――没有个不散的筵席。’}
}
{
    \eA{even had she not been so diligent and careful, she couldn't have been set
        aside! But what is provoking is that that lot,}
}
{
}

\Verse{38}
{
    \zA{谁守一辈子呢？}
}
{
    \eA{like Ch'ing Wen and Ch'i Hsia,}
}
{
}

\Verse{39}
{
    \zA{不过三年五载，各人干各人的去了，那时谁还管谁呢？”}
}
{
    \eA{should have been included in the upper class. Yet it's because every one
        places such reliance on the fine reputation of their father and mother that
        they exalt them.}
}
{
}

\Verse{40}
{
    \zA{这两句话 不觉感动了佳蕙心肠，由不得眼圈儿红了，又不好意思无端的哭，只得勉强笑道： “你这话说的是。}
}
{
    \eA{Now, do tell me, is this sufficient to anger one or not?" "It won't do to be
        angry with them!" Hsiao Hung observed. "The proverb says: 'You may erect a
        shed a thousand \_li\_ long, but there is no entertainment from which the
        guests will not disperse!'}
}
{
}

\Verse{41}
{
    \zA{昨儿宝玉还说：明儿怎么收拾房子，怎么做衣裳。}
}
{
    \eA{And who is it that will tarry here for a whole lifetime? In another three
        years or five years every single one of us will have gone her own way;}
}
{
}

\Verse{42}
{
    \zA{倒像有几百年 熬煎似的。”}
}
{
    \eA{and who will, when that time comes, worry her mind about any one else?"}
}
{
}

\Verse{43}
{
    \zA{小红听了，冷笑两声，方要说话，只见一个未留头的小丫头走进来，手里拿着 些花样子并两张纸，说道：“这两个花样子叫你描出来呢。”}
}
{
    \eA{These allusions had the unexpected effect of touching Chia Hui to the heart;
        and in spite of herself the very balls of her eyes got red. But so uneasy
        did she feel at crying for no reason that she had to exert herself to force
        a smile. "What you say is true," she ventured.}
}
{
}

\Verse{44}
{
    \zA{说着，向小红撂下， 回转身就跑了。}
}
{
    \eA{"And yet, Pao-yü even yesterday explained how the rooms should be arranged
        by and bye; and how the clothes should be made,}
}
{
}

\Verse{45}
{
    \zA{小红向外问道：“到底是谁的？}
}
{
    \eA{just as if he was bound to hang on to dear life for several hundreds of
        years."}
}
{
}

\Verse{46}
{
    \zA{也等不的说完就跑。}
}
{
    \eA{Hsiao Hung, at these words, gave a couple of sardonic smiles.}
}
{
}

\Verse{47}
{
    \zA{‘谁蒸下馒头 等着你――怕冷了不成？’”}
}
{
    \eA{But when about to pass some remark, she perceived a youthful servant-girl,
        who had not as yet let her hair grow,}
}
{
}

\Verse{48}
{
    \zA{那小丫头在窗外只说得一声：“是绮大姐姐的。”}
}
{
    \eA{walk in, holding in her hands several patterns and two sheets of paper. "You
        are asked," she said, "to trace these two designs!"}
}
{
}

\Verse{49}
{
    \zA{抬 起脚来，咕咚咕咚又跑了。}
}
{
    \eA{As she spoke, she threw them at Hsiao Hung, and twisting herself round,}
}
{
}

\Verse{50}
{
    \zA{小红便赌气把那样子撂在一边，向抽屉内找笔。}
}
{
    \eA{she immediately scampered away. "Whose are they, after all?" Hsiao Hung
        inquired, addressing herself outside.}
}
{
}

\Verse{51}
{
    \zA{找了半 天，都是秃的，因说道：“前儿一枝新笔放在那里了？}
}
{
    \eA{"Couldn't you wait even so much as to conclude what you had to say, but flew
        off at once? Who is steaming bread and waiting for you? Or are you afraid,
        forsooth,}
}
{
}

\Verse{52}
{
    \zA{怎么想不起来？”}
}
{
    \eA{lest it should get cold?" "They belong to sister Ch'i,"}
}
{
}

\Verse{53}
{
    \zA{一面说， 一面出神，想了一回，方笑道：“是了，前儿晚上莺儿拿了去了。”}
}
{
    \eA{the young servant-girl merely returned for answer from outside the window;
        and raising her feet high, she ran tramp-tramp on her way back again. Hsiao
        Hung lost control over her temper,}
}
{
}

\Verse{54}
{
    \zA{因向佳蕙道： “你替我取了来。”}
}
{
    \eA{and snatching the designs, she flung them on one side. She then rummaged in
        a drawer for a pencil,}
}
{
}

\Verse{55}
{
    \zA{佳蕙道：“花大姐姐还等着我替他拿箱子，你自己取去罢。”}
}
{
    \eA{but finding, after a prolonged search, that they were all blunt; "Where did
        I," she thereupon ejaculated, "put that brand-new pencil the other day?}
}
{
}

\Verse{56}
{
    \zA{小红道：“他等着你，你还坐着闲磕牙儿？}
}
{
    \eA{How is it I can't remember where it is?" While she soliloquised, she became
        wrapt in thought. After some reflection she,}
}
{
}

\Verse{57}
{
    \zA{我不叫你取去，他也不‘等’你了。}
}
{
    \eA{at length, gave a smile. "Of course!" she exclaimed, "the other evening Ying
        Erh took it away."}
}
{
}

\Verse{58}
{
    \zA{坏 透了的小蹄子！”}
}
{
    \eA{And turning towards Chia Hui, "Fetch it for me," she shouted.}
}
{
}

\Verse{59}
{
    \zA{说着自己便出房来。}
}
{
    \eA{"Sister Hua," Chia Hui rejoined, "is waiting for me to get a box for her,}
}
{
}

\Verse{60}
{
    \zA{出了怡红院，一径往宝钗院内来，刚至沁芳亭畔，只见宝 玉的奶娘李嬷嬷从那边来。}
}
{
    \eA{so you had better go for it yourself!" "What!" remarked Hsiao Hung, "she's
        waiting for you, and are you still squatting here chatting leisurely? Hadn't
        it been that I asked you to go and fetch it,}
}
{
}

\Verse{61}
{
    \zA{小红立住，笑问道：“李奶奶，你老人家那里去了？}
}
{
    \eA{she too wouldn't have been waiting for you; you most perverse vixen!" With
        these words on her lips,}
}
{
}

\Verse{62}
{
    \zA{怎 么打这里来？”}
}
{
    \eA{she herself walked out of the room, and leaving the I Hung court,}
}
{
}

\Verse{63}
{
    \zA{李嬷嬷站住，将手一拍，道：“你说，好好儿的，又看上了那个什 么‘云哥儿’‘雨哥儿’的，这会子逼着我叫了他来。}
}
{
    \eA{she straightway proceeded in the direction of Pao-ch'ai's court. As soon,
        however, as she reached the Hsin Fang pavilion, she saw dame Li, Pao-yü's
        nurse, appear in view from the opposite side; so Hsiao Hung halted and
        putting on a smile,}
}
{
}

\Verse{64}
{
    \zA{明儿叫上屋里听见，可又是 不好。”}
}
{
    \eA{"Nurse Li," she asked, "where are you, old dame, bound for? How is it you're
        coming this way?"}
}
{
}

\Verse{65}
{
    \zA{小红笑道：“你老人家当真的就信着他去叫么？”}
}
{
    \eA{Nurse Li stopped short, and clapped her hands. "Tell me," she said, "has he
        deliberately again gone and fallen in love with that Mr. something or other
        like Yun (cloud),}
}
{
}

\Verse{66}
{
    \zA{李嬷嬷道：“可怎么样 呢？”}
}
{
    \eA{or Yü (rain)? They now insist upon my bringing him inside,}
}
{
}

\Verse{67}
{
    \zA{小红笑道：“那一个要是知好歹，就不进来才是。”}
}
{
    \eA{but if they get wind of it by and bye in the upper rooms, it won't again be
        a nice thing." "Are you, old lady," replied Hsiao Hung smiling,}
}
{
}

\Verse{68}
{
    \zA{李嬷嬷道：“他又不傻， 为什么不进来？”}
}
{
    \eA{"taking things in such real earnest that you readily believe them and want
        to go and ask him in here?" "What can I do?" rejoined nurse Li.}
}
{
}

\Verse{69}
{
    \zA{小红道：“既是进来，你老人家该别和他一块儿来；回来叫他一 个人混碰，看他怎么样！”}
}
{
    \eA{"Why, that fellow," added Hsiao Hung laughingly, "will, if he has any idea
        of decency, do the right thing and not come." "Besides, he's not a fool!"
        pleaded nurse Li; "so why shouldn't he come in?" "Well, if he is to come,"}
}
{
}

\Verse{70}
{
    \zA{李嬷嬷道：“我有那样大工夫和他走!不过告诉了他， 回来打发个小丫头子，或是老婆子，带进他来就完了。”}
}
{
    \eA{answered Hsiao Hung, "it will devolve upon you, worthy dame, to lead him
        along with you; for were you by and bye to let him penetrate inside all
        alone and knock recklessly about, why, it won't do at all." "Have I got all
        that leisure,"}
}
{
}

\Verse{71}
{
    \zA{说着拄着拐一径去了。}
}
{
    \eA{retorted nurse Li, "to trudge along with him? I'll simply tell him to come;}
}
{
}

\Verse{72}
{
    \zA{小红听说，便站着出神，且不去取笔。}
}
{
    \eA{and later on I can despatch a young servant-girl or some old woman to bring
        him in, and have done." Saying this,}
}
{
}

\Verse{73}
{
    \zA{不多时，只见一个小丫头跑来，见小红 站在那里，便问道：“红姐姐，你在这里作什么呢？”}
}
{
    \eA{she continued her way, leaning on her staff. After listening to her
        rejoinder, Hsiao Hung stood still; and plunging in abstraction, she did not
        go and fetch the pencil. But presently, she caught sight of a servant-girl
        running that way.}
}
{
}

\Verse{74}
{
    \zA{小红抬头见是小丫头子坠儿， 小红道：“那里去？”}
}
{
    \eA{Espying Hsiao Hung lingering in that spot, "Sister Hung," she cried, "what
        are you doing in here?" Hsiao Hung raised her head,}
}
{
}

\Verse{75}
{
    \zA{坠儿道：“叫我带进芸二爷来。”}
}
{
    \eA{and recognised a young waiting-maid called Chui Erh. "Where are you off
        too?"}
}
{
}

\Verse{76}
{
    \zA{说着，一径跑了。}
}
{
    \eA{Hsiao Hung asked. "I've been told to bring in master Secundus,}
}
{
}

\Verse{77}
{
    \zA{这里小 红刚走至蜂腰桥门前，只见那边坠儿引着贾芸来了。}
}
{
    \eA{Mr. Yün," Chui Erh replied. After which answer, she there and then departed
        with all speed. Hsiao Hung reached, meanwhile,}
}
{
}

\Verse{78}
{
    \zA{那贾芸一面走，一面拿眼把小 红一溜；那小红只装着和坠儿说话，也把眼去一溜贾芸：四目恰好相对。}
}
{
    \eA{the Feng Yao bridge. As soon as she approached the gateway, she perceived
        Chui Erh coming along with Chia Yün from the opposite direction. While
        advancing Chia Yün ogled Hsiao Hung; and Hsiao Hung too, though pretending
        to be addressing herself to Chui Erh,}
}
{
}

\Verse{79}
{
    \zA{小红不觉 把脸一红，一扭身往蘅芜院去了。}
}
{
    \eA{cast a glance at Chia Yün; and their four eyes, as luck would have it, met.
        Hsiao Hung involuntarily blushed all over;}
}
{
}

\Verse{80}
{
    \zA{不在话下。}
}
{
    \eA{and turning herself round,}
}
{
}

\Verse{81}
{
    \zA{这里贾芸随着坠儿逶迤来至怡红院中，坠儿先进去回明了，然后方领贾芸进 去。}
}
{
    \eA{she walked off towards the Heng Wu court. But we will leave her there
        without further remarks. During this time, Chia Yün followed Chui Erh, by a
        circuitous way, into the I Hung court.}
}
{
}

\Verse{82}
{
    \zA{贾芸看时，只见院内略略有几点山石，种着芭蕉，那边有两只仙鹤，在松树下 剔翎。}
}
{
    \eA{Chui Erh entered first and made the necessary announcement. Then
        subsequently she ushered in Chia Yün. When Chia Yün scrutinised the
        surroundings,}
}
{
}

\Verse{83}
{
    \zA{一溜回廊上吊着各色笼子，笼着仙禽异鸟。}
}
{
    \eA{he perceived, here and there in the court, several blocks of rockery, among
        which were planted banana-trees.}
}
{
}

\Verse{84}
{
    \zA{上面小小五间抱厦，一色雕镂新 鲜花样？}
}
{
    \eA{On the opposite side were two storks preening their feathers under the fir
        trees. Under the covered passage were suspended,}
}
{
}

\Verse{85}
{
    \zA{扇，上面悬着一个匾，四个大字，题道是：“怡红快绿。”}
}
{
    \eA{in a row, cages of every description, containing all sorts of fairylike,
        rare birds. In the upper part were five diminutive anterooms,}
}
{
}

\Verse{86}
{
    \zA{贾芸想道：“怪 道叫‘怡红院’，原来匾上是这四个字。”}
}
{
    \eA{uniformly carved with, unique designs; and above the framework of the door
        was hung a tablet with the inscription in four huge characters—"I Hung K'uai
        Lü,}
}
{
}

\Verse{87}
{
    \zA{正想着，只听里面隔着纱窗子笑说道： “快进来罢，我怎么就忘了你两三个月！”}
}
{
    \eA{the happy red and joyful green." "I thought it strange," Chia Yün argued
        mentally, "that it should be called the I Hung court; but are these, in
        fact, the four characters inscribed on the tablet!"}
}
{
}

\Verse{88}
{
    \zA{贾芸听见是宝玉的声音，连忙进入房内， 抬头一看，只见金碧辉煌，文章？}
}
{
    \eA{But while he was communing within himself, he heard some one laugh and then
        exclaim from the inner side of the gauze window: "Come in at once! How is it
        that I've forgotten you these two or three months?"}
}
{
}

\Verse{89}
{
    \zA{烁，却看不见宝玉在那里。}
}
{
    \eA{As soon as Chia Yün recognised Pao-yü's voice, he entered the room with
        hurried step.}
}
{
}

\Verse{90}
{
    \zA{一回头，只见左边立 着一架大穿衣镜，从镜后转出两个一对儿十五六岁的丫头来，说：“请二爷里头屋 里坐。”}
}
{
    \eA{On raising his head, his eye was attracted by the brilliant splendour
        emitted by gold and jade and by the dazzling lustre of the elegant
        arrangements. He failed, however, to detect where Pao-yü was ensconced. The
        moment he turned his head round, he espied, on the left side,}
}
{
}

\Verse{91}
{
    \zA{贾芸连正眼也不敢看，连忙答应了。}
}
{
    \eA{a large cheval-glass; behind which appeared to view, standing side by side,
        two servant-girls of fifteen or sixteen years of age.}
}
{
}

\Verse{92}
{
    \zA{又进一道碧纱厨，只见小小一张填漆床上，悬着大红销金撒花帐子，宝玉穿着 家常衣服，？}
}
{
    \eA{"Master Secundus," they ventured, "please take a seat in the inner room."
        Chia Yün could not even muster courage to look at them straight in the face;
        but promptly assenting, he walked into a green gauze mosquito-house,}
}
{
}

\Verse{93}
{
    \zA{着鞋，倚在床上，拿着本书；看见他进来，将书掷下，早带笑立起身 来。}
}
{
    \eA{where he saw a small lacquered bed, hung with curtains of a deep red colour,
        with clusters of flowers embroidered in gold. Pao-yü, wearing a house-dress
        and slipshod shoes,}
}
{
}

\Verse{94}
{
    \zA{贾芸忙上前请了安，宝玉让坐，便在下面一张椅子上坐了。}
}
{
    \eA{was reclining on the bed, a book in hand. The moment he perceived Chia Yün
        walk in, he discarded his book, and forthwith smiled and raised himself up.}
}
{
}

\Verse{95}
{
    \zA{宝玉笑道：“只从 那个月见了你，我叫你往书房里来，谁知接接连连许多事情，就把你忘了。”}
}
{
    \eA{Chia Yün hurriedly pressed forward and paid his salutation. Pao-yü then
        offered him a seat; but he simply chose a chair in the lower part of the
        apartment. "Ever since the moon in which I came across you,"}
}
{
}

\Verse{96}
{
    \zA{贾芸 笑道：“总是我没造化，偏又遇着叔叔欠安。}
}
{
    \eA{Pao-yü observed smilingly, "and told you to come into the library, I've had,
        who would have thought it,}
}
{
}

\Verse{97}
{
    \zA{叔叔如今可大安了？”}
}
{
    \eA{endless things to continuously attend to, so that I forgot all about you."}
}
{
}

\Verse{98}
{
    \zA{宝玉道：“大 好了。}
}
{
    \eA{"It's I, indeed, who lacked good fortune!"}
}
{
}

\Verse{99}
{
    \zA{我倒听见说你辛苦了好几天。”}
}
{
    \eA{rejoined Chia Yün, with a laugh; "particularly so, as it again happened that
        you,}
}
{
}

\Verse{100}
{
    \zA{贾芸道：“辛苦也是该当的。}
}
{
    \eA{uncle, fell ill. But are you quite right once more?" "All right!" answered
        Pao-yü.}
}
{
}

\Verse{101}
{
    \zA{叔叔大安了， 也是我们一家子的造化。”}
}
{
    \eA{"I heard that you've been put to much trouble and inconvenience on a good
        number of days!" "Had I even had any trouble to bear,"}
}
{
}

\Verse{102}
{
    \zA{说着，只见有个丫鬟端了茶来与他。}
}
{
    \eA{added Chia Yün, "it would have been my duty to bear it. But your complete
        recovery,}
}
{
}

\Verse{103}
{
    \zA{那贾芸嘴里和宝玉 说话，眼睛却瞅那丫鬟：细挑身子，容长脸儿，穿着银红袄儿，青缎子坎肩，白绫 细褶儿裙子。}
}
{
    \eA{uncle, is really a blessing to our whole family." As he spoke, he discerned
        a couple of servant-maids come to help him to a cup of tea. But while
        conversing with Pao-yü, Chia Yün was intent upon scrutinising the girl with
        slim figure, and oval face, and clad in a silvery-red jacket,}
}
{
}

\Verse{104}
{
    \zA{那贾芸自从宝玉病了，他在里头混了两天，都把有名人口记了一半， 他看见这丫鬟，知道是袭人。}
}
{
    \eA{a blue satin waistcoat and a white silk petticoat with narrow pleats. At the
        time of Pao-yü's illness, Chia Yün had spent a couple of days in the inner
        apartments, so that he remembered half of the inmates of note, and the
        moment he set eyes upon this servant-girl he knew that it was Hsi Jen;}
}
{
}

\Verse{105}
{
    \zA{他在宝玉房中比别人不同，如今端了茶来，宝玉又在 旁边坐着，便忙站起来笑道：“姐姐怎么给我倒起茶来？}
}
{
    \eA{and that she was in Pao-yü's rooms on a different standing to the rest. Now
        therefore that she brought the tea in herself and that Pao-yü was, besides,
        sitting by, he rose to his feet with alacrity and put on a smile. "Sister,"
        he said, "how is it that you are pouring tea for me?}
}
{
}

\Verse{106}
{
    \zA{我来到叔叔这里，又不是 客，等我自己倒罢了。”}
}
{
    \eA{I came here to pay uncle a visit; what's more I'm no stranger, so let me
        pour it with my own hands!" "Just you sit down and finish!"}
}
{
}

\Verse{107}
{
    \zA{宝玉道：“你只管坐着罢。}
}
{
    \eA{Pao-yü interposed; "will you also behave in this fashion with
        servant-girls?"}
}
{
}

\Verse{108}
{
    \zA{丫头们跟前也是这么着。”}
}
{
    \eA{"In spite of what you say;" remarked Chia Yün smiling,}
}
{
}

\Verse{109}
{
    \zA{贾 芸笑道：“虽那么说，叔叔屋里的姐姐们，我怎么敢放肆呢。”}
}
{
    \eA{"they are young ladies attached to your rooms, uncle, and how could I
        presume to be disorderly in my conduct?" So saying, he took a seat and drank
        his tea.}
}
{
}

\Verse{110}
{
    \zA{一面说，一面坐下 吃茶。}
}
{
    \eA{Pao-yü then talked to him about trivial and irrelevant matters;}
}
{
}

\Verse{111}
{
    \zA{那宝玉便和他说些没要紧的散话：又说道谁家的戏子好，谁家的花园好，又告 诉他谁家的丫头标致，谁家的酒席丰盛，又是谁家有奇货，又是谁家有异物。}
}
{
    \eA{and afterwards went on to tell him in whose household the actresses were
        best, and whose gardens were pretty. He further mentioned to him in whose
        quarters the servant-girls were handsome, whose banquets were sumptuous,}
}
{
}

\Verse{112}
{
    \zA{那贾 芸口里只得顺着他说。}
}
{
    \eA{as well as in whose home were to be found strange things, and what family
        possessed remarkable objects.}
}
{
}

\Verse{113}
{
    \zA{说了一回，见宝玉有些懒懒的了，便起身告辞。}
}
{
    \eA{Chia Yün was constrained to humour him in his conversation; but after a
        chat, which lasted for some time,}
}
{
}

\Verse{114}
{
    \zA{宝玉也不甚 留，只说：“你明儿闲了只管来。”}
}
{
    \eA{he noticed that Pao-yü was somewhat listless, and he promptly stood up and
        took his leave. And Pao-yü too did not use much pressure to detain him.}
}
{
}

\Verse{115}
{
    \zA{仍命小丫头子坠儿送出去了。}
}
{
    \eA{"To-morrow, if you have nothing to do, do come over!" he merely observed;}
}
{
}

\Verse{116}
{
    \zA{贾芸出了怡红院，见四顾无人，便慢慢的停着些走，口里一长一短和坠儿说话。}
}
{
    \eA{after which, he again bade the young waiting-maid, Chui Erh, see him out.
        Having left the I Hung court, Chia Yün cast a glance all round; and,
        realising that there was no one about,}
}
{
}

\Verse{117}
{
    \zA{先问他：“几岁了？}
}
{
    \eA{he slackened his pace at once, and while proceeding leisurely,}
}
{
}

\Verse{118}
{
    \zA{名字叫什么？}
}
{
    \eA{he conversed, in a friendly way,}
}
{
}

\Verse{119}
{
    \zA{你父母在那行上？}
}
{
    \eA{with Chui Erh on one thing and another.}
}
{
}

\Verse{120}
{
    \zA{在宝叔屋里几年了？}
}
{
    \eA{First and foremost he inquired of her what was her age;}
}
{
}

\Verse{121}
{
    \zA{一个月多少钱？}
}
{
    \eA{and her name. "Of what standing are your father and mother?"}
}
{
}

\Verse{122}
{
    \zA{共总宝叔屋内有几个女孩子？”}
}
{
    \eA{he said, "How many years have you been in uncle Pao's apartments? How much
        money do you get a month?}
}
{
}

\Verse{123}
{
    \zA{那坠儿见问，便一桩桩的都告诉他了。}
}
{
    \eA{In all how many girls are there in uncle Pao's rooms?" As Chui Erh heard the
        questions set to her,}
}
{
}

\Verse{124}
{
    \zA{贾芸又道： “刚才那个和你说话的，他可是叫小红？”}
}
{
    \eA{she readily made suitable reply to each. "The one, who was a while back
        talking to you," continued Chia Yün, "is called Hsiao Hung,}
}
{
}

\Verse{125}
{
    \zA{坠儿笑道：“他就叫小红。}
}
{
    \eA{isn't she?" "Yes, her name is Hsiao Hung!" replied Chui Erh smiling;}
}
{
}

\Verse{126}
{
    \zA{你问他作什 么？”}
}
{
    \eA{"but why do you ask about her?"}
}
{
}

\Verse{127}
{
    \zA{贾芸道：“方才他问你什么绢子，我倒拣了一块。”}
}
{
    \eA{"She inquired of you just now about some handkerchief or other," answered
        Chia Yün; "well, I've picked one up." Chui Erh greeted this response with a
        smile.}
}
{
}

\Verse{128}
{
    \zA{坠儿听了笑道：“他问 了我好几遍：可有看见他的绢子的。}
}
{
    \eA{"Many are the times," she said; "that she has asked me whether I had seen
        her handkerchief; but have I got all that leisure to worry my mind about
        such things?}
}
{
}

\Verse{129}
{
    \zA{我那里那么大工夫管这些事？}
}
{
    \eA{She spoke to me about it again to-day; and she suggested that I should find
        it for her,}
}
{
}

\Verse{130}
{
    \zA{今儿他又问我， 他说我替他找着了他还谢我呢。}
}
{
    \eA{and that she would also recompense me. This she told me when we were just
        now at the entrance of the Heng Wu court, and you too, Mr. Secundus,}
}
{
}

\Verse{131}
{
    \zA{才在蘅芜院门口儿说的，二爷也听见了，不是我撒 谎。}
}
{
    \eA{overheard her, so that I'm not lying. But, dear Mr. Secundus, since you've
        picked it up, give it to me. Do! And I'll see what she will give me as a
        reward."}
}
{
}

\Verse{132}
{
    \zA{好二爷，你既拣了，给我罢，我看他拿什么谢我。”}
}
{
    \eA{The truth is that Chia Yün had, the previous moon when he had come into the
        garden to attend to the planting of trees, picked up a handkerchief,}
}
{
}

\Verse{133}
{
    \zA{原来上月贾芸进来种树之 时，便拣了一块罗帕，知是这园内的人失落的，但不知是那一个人的，故不敢造次。}
}
{
    \eA{which he conjectured must have been dropped by some inmate of those grounds;
        but as he was not aware whose it was, he did not consequently presume to act
        with indiscretion. But on this occasion, he overheard Hsiao Hung make
        inquiries of Chui Erh on the subject;}
}
{
}

\Verse{134}
{
    \zA{今听见小红问坠儿，知是他的，心内不胜喜幸。}
}
{
    \eA{and concluding that it must belong to her, he felt immeasurably delighted.
        Seeing, besides, how importunate Chui Erh was,}
}
{
}

\Verse{135}
{
    \zA{又见坠儿追索，心中早得了主意， 便向袖内将自己的一块取出来，向坠儿笑道：“我给是给你，你要得了他的谢礼， 可不许瞒着我。”}
}
{
    \eA{he at once devised a plan within himself, and vehemently producing from his
        sleeve a handkerchief of his own, he observed, as he turned towards Chui Erh
        with a smile: "As for giving it to you, I'll do so; but in the event of your
        obtaining any present from her,}
}
{
}

\Verse{136}
{
    \zA{坠儿满口里答应了，接了绢子，送出贾芸，回来找小红，不在话 下。}
}
{
    \eA{you mustn't impose upon me." Chui Erh assented to his proposal most
        profusely; and, taking the handkerchief, she saw Chia Yün out and then came
        back in search of Hsiao Hung.}
}
{
}

\Verse{137}
{
    \zA{如今且说宝玉打发贾芸去后，意思懒懒的，歪在床上，似有朦胧之态。}
}
{
    \eA{But we will leave her there for the present. We will now return to Pao-yü.
        After dismissing Chia Yün, he lay in such complete listlessness on the bed
        that he betrayed every sign of being half asleep.}
}
{
}

\Verse{138}
{
    \zA{袭人便 走上来，坐在床沿上推他，说道：“怎么又要睡觉？}
}
{
    \eA{Hsi Jen walked up to him, and seated herself on the edge of the bed, and
        pushing him, "What are you about to go to sleep again,"}
}
{
}

\Verse{139}
{
    \zA{你闷的很，出去逛逛不好？”}
}
{
    \eA{she said. "Would it not do your languid spirits good if you went out for a
        bit of a stroll?"}
}
{
}

\Verse{140}
{
    \zA{宝玉见说，携着他的手笑道：“我要去，只是舍不得你。”}
}
{
    \eA{Upon hearing her voice, Pao-yü grasped her hand in his. "I would like to go
        out," he smiled, "but I can't reconcile myself to the separation from you!"}
}
{
}

\Verse{141}
{
    \zA{袭人笑道：“你没别的 说了！”}
}
{
    \eA{"Get up at once!" laughed Hsi Jen. And as she uttered these words, she
        pulled Pao-yü up. "Where can I go?"}
}
{
}

\Verse{142}
{
    \zA{一面说，一面拉起他来。}
}
{
    \eA{exclaimed Pao-yü. "I'm quite surfeited with everything."}
}
{
}

\Verse{143}
{
    \zA{宝玉道：“可往那里去呢？}
}
{
    \eA{"Once out you'll be all right," Hsi Jen answered, "but if you simply give
        way to this languor,}
}
{
}

\Verse{144}
{
    \zA{怪腻腻烦烦的。”}
}
{
    \eA{you'll be more than ever sick of everything at heart."}
}
{
}

\Verse{145}
{
    \zA{袭 人道：“你出去了就好了。}
}
{
    \eA{Pao-yü could not do otherwise, dull and out of sorts though he was,}
}
{
}

\Verse{146}
{
    \zA{只管这么委琐，越发心里腻烦了。”}
}
{
    \eA{than accede to her importunities. Strolling leisurely out of the door of the
        room,}
}
{
}

\Verse{147}
{
    \zA{宝玉无精打彩，只 得依他。}
}
{
    \eA{he amused himself a little with the birds suspended under the verandah;}
}
{
}

\Verse{148}
{
    \zA{晃出了房门，在回廊上调弄了一回雀儿，出至院外，顺着沁芳溪，看了一 回金鱼。}
}
{
    \eA{then he wended his steps outside the court, and followed the course of the
        Hsin Fang stream; but after admiring the golden fish for a time, he espied,
        on the opposite hillock, two young deer come rushing down as swift as an
        arrow.}
}
{
}

\Verse{149}
{
    \zA{只见那边山坡上两只小鹿儿箭也似的跑来。}
}
{
    \eA{What they were up to Pao-yü could not discern; but while abandoning himself
        to melancholy, he caught sight of Chia Lan,}
}
{
}

\Verse{150}
{
    \zA{宝玉不解何意，正自纳闷，只 见贾兰在后面，拿着一张小弓儿赶来。}
}
{
    \eA{following behind, with a small bow in his hand, and hurrying down hill in
        pursuit of them. As soon as he realised that Pao-yü stood ahead of him, he
        speedily halted. "Uncle Secundus,"}
}
{
}

\Verse{151}
{
    \zA{一见宝玉在前，便站住了，笑道：“二叔叔 在家里呢，我只当出门去了呢。”}
}
{
    \eA{he smiled, "are you at home? I imagined you had gone out of doors!" "You are
        up to mischief again, eh?" Pao-yü rejoined. "They've done nothing to you,
        and why shoot at them with your arrows?"}
}
{
}

\Verse{152}
{
    \zA{宝玉道：“你又淘气了。}
}
{
    \eA{"I had no studies to attend to just now, so, being free with nothing to do,"}
}
{
}

\Verse{153}
{
    \zA{好好儿的，射他做什么？”}
}
{
    \eA{Chia Lan replied laughingly, "I was practising riding and archery."}
}
{
}

\Verse{154}
{
    \zA{贾兰笑道：“这会子不念书，闲着做什么？}
}
{
    \eA{"Shut up!" exclaimed Pao-yü. "When are you not engaged in practising?"
        Saying this, he continued his way and straightway reached the entrance of a
        court.}
}
{
}

\Verse{155}
{
    \zA{所以演习演习骑射。”}
}
{
    \eA{Here the bamboo foliage was thick, and the breeze sighed gently.}
}
{
}

\Verse{156}
{
    \zA{宝玉道：“磕了 牙，那时候儿才不演呢。”}
}
{
    \eA{This was the Hsiao Hsiang lodge. Pao-yü listlessly rambled in. He saw a
        bamboo portière hanging down to the ground.}
}
{
}

\Verse{157}
{
    \zA{说着，便顺脚一径来至一个院门前，看那凤尾森森，龙吟细细：正是潇湘馆。}
}
{
    \eA{Stillness prevailed. Not a human voice fell on the ear. He advanced as far
        as the window. Noticing that a whiff of subtle scent stole softly through
        the green gauze casement, Pao-yü applied his face closely against the frame
        to peep in,}
}
{
}

\Verse{158}
{
    \zA{宝玉信步走入，只见湘帘垂地，悄无人声。}
}
{
    \eA{but suddenly he caught the faint sound of a deep sigh and the words: "Day
        after day my feelings slumber drowsily!"}
}
{
}

\Verse{159}
{
    \zA{走至窗前，觉得一缕幽香从碧纱窗中暗 暗透出，宝玉便将脸贴在纱窗上。}
}
{
    \eA{Upon overhearing this exclamation, Pao-yü unconsciously began to feel a prey
        to inward longings; but casting a second glance, he saw Tai-yü stretching
        herself on the bed.}
}
{
}

\Verse{160}
{
    \zA{看时，耳内忽听得细细的长叹了一声，道：“‘每 日家情思睡昏昏！’”}
}
{
    \eA{"Why is it," smiled Pao-yü, from outside the window, "that your feelings day
        after day slumber drowsily?" So saying, he raised the portière and stepped
        in.}
}
{
}

\Verse{161}
{
    \zA{宝玉听了，不觉心内痒将起来。}
}
{
    \eA{The consciousness that she had not been reticent about her feelings made
        Tai-yü unwittingly flush scarlet.}
}
{
}

\Verse{162}
{
    \zA{再看时，只见黛玉在床上伸 懒腰。}
}
{
    \eA{Taking hold of her sleeve, she screened her face; and, turning her body
        round towards the inside,}
}
{
}

\Verse{163}
{
    \zA{宝玉在窗外笑道：“为什么‘每日家情思睡昏昏’的？”}
}
{
    \eA{she pretended to be fast asleep. Pao-yü drew near her. He was about to pull
        her round when he saw Tai-yü's nurse enter the apartment,}
}
{
}

\Verse{164}
{
    \zA{一面说，一面掀帘 子进来了。}
}
{
    \eA{followed by two matrons. "Is Miss asleep?" they said. "If so, we'll ask her
        over,}
}
{
}

\Verse{165}
{
    \zA{黛玉自觉忘情，不觉红了脸，拿袖子遮了脸，翻身向里装睡着了。}
}
{
    \eA{when she wakes up." As these words were being spoken, Tai-yü eagerly twisted
        herself round and sat up. "Who's asleep?" she laughed. "We thought you were
        fast asleep,}
}
{
}

\Verse{166}
{
    \zA{宝玉 才走上来，要扳他的身子，只见黛玉的奶娘并两个婆子却跟进来了，说：“妹妹睡 觉呢，等醒来再请罢。”}
}
{
    \eA{Miss," smiled the two or three matrons as soon as they perceived Tai-yü get
        up. This greeting over, they called Tzu Chüan. "Your young mistress," they
        said, "has awoke; come in and wait on her!"}
}
{
}

\Verse{167}
{
    \zA{刚说着，黛玉便翻身坐起来，笑道：“谁睡觉呢？”}
}
{
    \eA{While calling her, they quitted the room in a body. Tai-yü remained seated
        on the bed. Raising her arms, she adjusted her hair,}
}
{
}

\Verse{168}
{
    \zA{那两 三个婆子见黛玉起来，便笑道：“我们只当姑娘睡着了。”}
}
{
    \eA{and smilingly she observed to Pao-yü, "When people are asleep, what do you
        walk in for?" At the sight of her half-closed starlike eyes and of her
        fragrant cheeks,}
}
{
}

\Verse{169}
{
    \zA{说着，便叫紫鹃说：“姑 娘醒了，进来伺候。”}
}
{
    \eA{suffused with a crimson blush, Pao-yü's feelings were of a sudden awakened;
        so, bending his body, he took a seat on a chair,}
}
{
}

\Verse{170}
{
    \zA{一面说，一面都去了。}
}
{
    \eA{and asked with a smile: "What were you saying a short while back?"}
}
{
}

\Verse{171}
{
    \zA{黛玉坐在床上，一面抬手整理鬓发，一面笑向宝玉道：“人家睡觉，你进来做 什么？”}
}
{
    \eA{"I wasn't saying anything," Tai-yü replied. "What a lie you're trying to ram
        down my throat!" laughed Pao-yü. "I heard all." But in the middle of their
        colloquy, they saw Tzu Chüan enter.}
}
{
}

\Verse{172}
{
    \zA{宝玉见他星眼微饧，香腮带赤，不觉神魂早荡，一歪身坐在椅子上，笑道： “你才说什么？”}
}
{
    \eA{Pao-yü then put on a smiling face. "Tzu Chüan!" he cried, "pour me a cup of
        your good tea!" "Where's the good tea to be had?" Tzu Chüan answered. "If
        you want good tea, you'd better wait till Hsi Jen comes."}
}
{
}

\Verse{173}
{
    \zA{黛玉道：“我没说什么。”}
}
{
    \eA{"Don't heed him!" interposed Tai-yü. "Just go first and draw me some water."}
}
{
}

\Verse{174}
{
    \zA{宝玉笑道：“给你个榧子吃呢!我都 听见了。”}
}
{
    \eA{"He's a visitor," remonstrated Tzu Chüan, "and, of course, I should first
        pour him a cup of tea, and then go and draw the water."}
}
{
}

\Verse{175}
{
    \zA{二人正说话，只见紫鹃进来，宝玉笑道：“紫鹃，把你们的好茶沏碗我 喝。”}
}
{
    \eA{With this answer, she started to serve the tea. "My dear girl," Pao-yü
        exclaimed laughingly, "If I could only share the same bridal curtain with
        your lovable young mistress, would I ever be able (to treat you as a
        servant) by making you fold the covers and make the beds."}
}
{
}

\Verse{176}
{
    \zA{紫鹃道：“我们那里有好的？}
}
{
    \eA{Lin Tai-yü at once drooped her head. "What are you saying?" she
        remonstrated.}
}
{
}

\Verse{177}
{
    \zA{要好的只好等袭人来。”}
}
{
    \eA{"What, did I say anything?" smiled Pao-yü. Tai-yü burst into tears.}
}
{
}

\Verse{178}
{
    \zA{黛玉道：“别理他。}
}
{
    \eA{"You've recently," she observed,}
}
{
}

\Verse{179}
{
    \zA{你先给我舀水去罢。”}
}
{
    \eA{"got into a new way. Whatever slang you happen to hear outside you come and
        tell me.}
}
{
}

\Verse{180}
{
    \zA{紫鹃道：“他是客，自然先沏了茶来再舀水去。”}
}
{
    \eA{And whenever you read any improper book, you poke your fun at me. What! have
        I become a laughing-stock for gentlemen!"}
}
{
}

\Verse{181}
{
    \zA{说着，倒 茶去了。}
}
{
    \eA{As she began to cry, she jumped down from bed,}
}
{
}

\Verse{182}
{
    \zA{宝玉笑道：“好丫头!‘若共你多情小姐同鸳帐，怎舍得叫你叠被铺床？’”}
}
{
    \eA{and promptly left the room. Pao-yü was at a loss how to act. So agitated was
        he that he hastily ran up to her, "My dear cousin," he pleaded, "I do
        deserve death; but don't go and tell any one!}
}
{
}

\Verse{183}
{
    \zA{黛玉登时急了，撂下脸来说道：“你说什么？”}
}
{
    \eA{If again I venture to utter such kind of language, may blisters grow on my
        mouth and may my tongue waste away!"}
}
{
}

\Verse{184}
{
    \zA{宝玉笑道：“我何尝说什么？”}
}
{
    \eA{But while appealing to her feelings, he saw Hsi Jen approach him.}
}
{
}

\Verse{185}
{
    \zA{黛 玉便哭道：“如今新兴的，外头听了村话来，也说给我听；看了混帐书，也拿我取 笑儿。}
}
{
    \eA{"Go back at once," she cried, "and put on your clothes as master wants to
        see you." At the very mention of his father, Pao-yü felt suddenly as if
        struck by lightning. Regardless of everything and anything,}
}
{
}

\Verse{186}
{
    \zA{我成了替爷们解闷儿的了。”}
}
{
    \eA{he rushed, as fast as possible, back to his room, and changing his clothes,}
}
{
}

\Verse{187}
{
    \zA{一面哭，一面下床来，往外就走。}
}
{
    \eA{he came out into the garden. Here he discovered Pei Ming, standing at the
        second gateway,}
}
{
}

\Verse{188}
{
    \zA{宝玉心下慌 了，忙赶上来说：“好妹妹，我一时该死，你好歹别告诉去!我再敢说这些话，嘴 上就长个疔，烂了舌头。”}
}
{
    \eA{waiting for him. "Do you perchance know what he wants me for?" Pao-yü
        inquired. "Master, hurry out at once!" Pei Ming replied. "You must, of
        course, go and see him. When you get there, you are sure to find out what
        it's all about."}
}
{
}

\Verse{189}
{
    \zA{正说着，只见袭人走来，说道：“快回去穿衣裳去罢，老爷叫你呢。”}
}
{
    \eA{This said, he urged Pao-yü on, and together they turned past the large
        pavilion. Pao-yü was, however, still labouring under suspicion, when he
        heard, from the corner of the wall,}
}
{
}

\Verse{190}
{
    \zA{宝玉听 了，不觉打了个焦雷一般，也顾不得别的，疾忙回来穿衣服。}
}
{
    \eA{a loud outburst of laughter. Upon turning his head round, he caught sight of
        Hsüeh P'an jump out, clapping his hands. "Hadn't I said that my uncle wanted
        you?"}
}
{
}

\Verse{191}
{
    \zA{出园来，只见焙茗在 二门前等着。}
}
{
    \eA{he laughed. "Would you ever have rushed out with such alacrity?"}
}
{
}

\Verse{192}
{
    \zA{宝玉问道：“你可知道老爷叫我是为什么？”}
}
{
    \eA{Pei Ming also laughed, and fell on his knees. But Pao-yü remained for a long
        time under the spell of utter astonishment,}
}
{
}

\Verse{193}
{
    \zA{焙茗道：“爷快出来罢， 横竖是见去的，到那里就知道了。”}
}
{
    \eA{before he, at length, realised that it was Hsüeh P'au who had inveigled him
        to come out. Hsüeh P'an hastily made a salutation and a curtsey, and
        confessed his fault.}
}
{
}

\Verse{194}
{
    \zA{一面说，一面催着宝玉。}
}
{
    \eA{He next gave way to entreaties, saying: "Don't punish the young servant,}
}
{
}

\Verse{195}
{
    \zA{转过大厅，宝玉心里 还自狐疑，只听墙角边一阵呵呵大笑，回头见薛蟠拍着手跳出来，笑道：“要不说 姨夫叫你，你那里肯出来的这么快！”}
}
{
    \eA{for it is simply I who begged him go." Pao-yü too had then no other
        alternative but to smile. "I don't mind your playing your larks on me; but
        why," he inquired, "did you mention my father? Were I to go and tell my
        aunt, your mother, to see to the rights and the wrongs of the case,}
}
{
}

\Verse{196}
{
    \zA{焙茗也笑着跪下了。}
}
{
    \eA{how would you like it?" "My dear cousin," remarked Hsüeh P'an vehemently,}
}
{
}

\Verse{197}
{
    \zA{宝玉怔了半天，方想过 来，是薛蟠哄出他来。}
}
{
    \eA{"the primary idea I had in view was to ask you to come out a moment sooner
        and I forgot to respectfully shun the expression.}
}
{
}

\Verse{198}
{
    \zA{薛蟠连忙打恭作揖赔不是，又求：“别难为了小子，都是我 央及他去的。”}
}
{
    \eA{But by and bye, when you wish to chaff me, just you likewise allude to my
        father, and we'll thus be square." "Ai-ya!" exclaimed Pao-yü. "You do more
        than ever deserve death!!"}
}
{
}

\Verse{199}
{
    \zA{宝玉也无法了，只好笑问道：“你哄我也罢了，怎么说是老爷呢？}
}
{
    \eA{Then turning again towards Pei Ming, "You ruffian!" he said, "what are you
        still kneeling for?" Pei Ming began to bump his head on the ground with
        vehemence.}
}
{
}

\Verse{200}
{
    \zA{我告诉姨娘去，评评这个理，可使得么？”}
}
{
    \eA{"Had it been for anything else," Hsüeh P'an chimed in, "I wouldn't have made
        bold to disturb you;}
}
{
}

\Verse{201}
{
    \zA{薛蟠忙道：“好兄弟，我原为求你快些 出来，就忘了忌讳这句话，改日你要哄我，也说我父亲，就完了。”}
}
{
    \eA{but it's simply in connection with my birthday which is to-morrow, the third
        day of the fifth moon. Ch'eng Jih-hsing, who is in that curio shop of ours,
        unexpectedly brought along, goodness knows where he fished them from, fresh
        lotus so thick and so long,}
}
{
}

\Verse{202}
{
    \zA{宝玉道：“嗳 哟，越发的该死了。”}
}
{
    \eA{so mealy and so crisp; melons of this size; and a Siamese porpoise, that
        long and that big,}
}
{
}

\Verse{203}
{
    \zA{又向焙茗道：“反叛杂种，还跪着做什么？”}
}
{
    \eA{smoked with cedar, such as is sent as tribute from the kingdom of Siam. Are
        not these four presents,}
}
{
}

\Verse{204}
{
    \zA{焙茗连忙叩头 起来。}
}
{
    \eA{pray, rare delicacies? The porpoise is not only expensive,}
}
{
}

\Verse{205}
{
    \zA{薛蟠道：“要不是，我也不敢惊动：只因明儿五月初三日，是我的生日，谁知 老胡和老程他们，不知那里寻了来的：这么粗这么长粉脆的鲜藕，这么大的西瓜，
        这么长这么大的暹罗国进贡的灵柏香熏的暹罗猪、鱼。}
}
{
    \eA{but difficult to get, and that kind of lotus and melon must have cost him no
        end of trouble to grow! I lost no time in presenting some to my mother, and
        at once sent some to your old grandmother, and my aunt. But a good many of
        them still remain now; and were I to eat them all alone,}
}
{
}

\Verse{206}
{
    \zA{你说这四样礼物，可难得不 难得？}
}
{
    \eA{it would, I fear, be more than I deserve; so I concluded, after thinking
        right and left,}
}
{
}

\Verse{207}
{
    \zA{那鱼、猪不过贵而难得，这藕和瓜亏他怎么种出来的!我先孝敬了母亲，赶着 就给你们老太太、姨母送了些去。}
}
{
    \eA{that there was, besides myself, only you good enough to partake of some.
        That is why I specially invite you to taste them. But, as luck would have
        it, a young singing-boy has also come, so what do you say to you and I
        having a jolly day of it?"}
}
{
}

\Verse{208}
{
    \zA{如今留了些，我要自己吃恐怕折福，左思右想除 我之外惟你还配吃。}
}
{
    \eA{As they talked, they walked; and, as they walked, they reached the interior
        of the library. Here they discovered a whole assemblage consisting of Tan
        Kuang, Ch'eng Jih-hsing, Hu Ch'i-lai,}
}
{
}

\Verse{209}
{
    \zA{所以特请你来。}
}
{
    \eA{Tan T'ing-jen and others, and the singing-boy as well.}
}
{
}

\Verse{210}
{
    \zA{可巧唱曲儿的一个小子又来了，我和你乐一天 何如？”}
}
{
    \eA{As soon as these saw Pao-yü walk in, some paid their respects to him; others
        inquired how he was; and after the interchange of salutations,}
}
{
}

\Verse{211}
{
    \zA{一面说，一面来到他书房里，只见詹光、程日兴、胡斯来、单聘仁等并唱曲儿 的小子都在这里。}
}
{
    \eA{tea was drunk. Hsüeh P'an then gave orders to serve the wine. Scarcely were
        the words out of his mouth than the servant-lads bustled and fussed for a
        long while laying the table. When at last the necessary arrangements had
        been completed, the company took their seats. Pao-yü verily found the melons
        and lotus of an exceptional description.}
}
{
}

\Verse{212}
{
    \zA{见他进来，请安的，问好的，都彼此见过了。}
}
{
    \eA{"My birthday presents have not as yet been sent round," he felt impelled to
        say, a smile on his lips, "and here I come,}
}
{
}

\Verse{213}
{
    \zA{吃了茶，薛蟠即命 人：“摆酒来。”}
}
{
    \eA{ahead of them, to trespass on your hospitality." "Just so!" retorted Hsüeh
        P'an,}
}
{
}

\Verse{214}
{
    \zA{话犹未了，众小厮七手八脚摆了半天，方才停当归坐。}
}
{
    \eA{"but when you come to-morrow to congratulate me we'll consider what novel
        kind of present you can give me." "I've got nothing that I can give you,"}
}
{
}

\Verse{215}
{
    \zA{宝玉果见 瓜藕新异，因笑道：“我的寿礼还没送来，倒先扰了。”}
}
{
    \eA{rejoined Pao-yü. "As far as money, clothes, eatables and other such articles
        go, they are not really mine: all I can call my own are such pages of
        characters that I may write,}
}
{
}

\Verse{216}
{
    \zA{薛蟠道：“可是呢，你明 儿来拜寿，打算送什么新鲜物儿？”}
}
{
    \eA{or pictures that I may draw." "Your reference to pictures," added Hsüeh P'an
        smiling, "reminds me of a book I saw yesterday,}
}
{
}

\Verse{217}
{
    \zA{宝玉道：“我没有什么送的。}
}
{
    \eA{containing immodest drawings; they were, truly, beautifully done.}
}
{
}

\Verse{218}
{
    \zA{若论银钱吃穿等 类的东西，究竟还不是我的；惟有写一张字，或画一张画，这才是我的。”}
}
{
    \eA{On the front page there figured also a whole lot of characters. But I didn't
        carefully look at them; I simply noticed the name of the person, who had
        executed them. It was, in fact, something or other like Keng Huang. The
        pictures were, actually,}
}
{
}

\Verse{219}
{
    \zA{薛蟠笑 道：“你提画儿，我才想起来了：昨儿我看见人家一本春宫儿，画的很好。}
}
{
    \eA{exceedingly good!" This allusion made Pao-yü exercise his mind with
        innumerable conjectures. "Of pictures drawn from past years to the present,
        I have," he said, "seen a good many, but I've never come across any Keng
        Huang."}
}
{
}

\Verse{220}
{
    \zA{上头还 有许多的字，我也没细看，只看落的款，原来是什么‘庚黄’的。}
}
{
    \eA{After considerable thought, he could not repress himself from bursting out
        laughing. Then asking a servant to fetch him a pencil, he wrote a couple of
        words on the palm of his hand.}
}
{
}

\Verse{221}
{
    \zA{真好的了不得。”}
}
{
    \eA{This done, he went on to inquire of Hsüeh.}
}
{
}

\Verse{222}
{
    \zA{宝玉听说，心下猜疑道：“古今字画也都见过些，那里有个‘庚黄’？”}
}
{
    \eA{P'an: "Did you see correctly that it read Keng Huang?" "How could I not have
        seen correctly?" ejaculated Hsüeh P'an. Pao-yü thereupon unclenched his hand
        and allowed him to peruse,}
}
{
}

\Verse{223}
{
    \zA{想了半天， 不觉笑将起来，命人取过笔来，在手心里写了两个字，又问薛蟠道：“你看真了是 ‘庚黄’么？”}
}
{
    \eA{what was written in it. "Were they possibly these two characters?" he
        remarked. "These are, in point of fact, not very dissimilar from what Keng
        Huang look like?" On scrutinising them, the company noticed the two words
        T'ang Yin, and they all laughed. "They must,}
}
{
}

\Verse{224}
{
    \zA{薛蟠道：“怎么没看真？”}
}
{
    \eA{we fancy, have been these two characters!" they cried. "Your eyes, Sir,}
}
{
}

\Verse{225}
{
    \zA{宝玉将手一撒给他看道：“可是这两个 字罢？}
}
{
    \eA{may, there's no saying, have suddenly grown dim!" Hsüeh P'an felt utterly
        abashed. "Who could have said," he smiled,}
}
{
}

\Verse{226}
{
    \zA{其实和‘庚黄’相去不远。”}
}
{
    \eA{"whether they were T'ang Yin or Kuo Yin, (candied silver or fruit silver)."
        As he cracked this joke,}
}
{
}

\Verse{227}
{
    \zA{众人都看时，原来是“唐寅”两个字，都笑道： “想必是这两个字，大爷一时眼花了，也未可知。”}
}
{
    \eA{however, a young page came and announced that Mr. Feng had arrived. Pao-yü
        concluded that the new comer must be Feng Tzu-ying, the son of Feng T'ang,
        general with the prefix of Shen Wu." "Ask him in at once," Hsüeh P'an and
        his companions shouted with one voice.}
}
{
}

\Verse{228}
{
    \zA{薛蟠自觉没趣，笑道：“谁知 他是‘糖银’是‘果银’的！”}
}
{
    \eA{But barely were these words out of their mouths, than they realised that
        Feng Tzu-ying had already stepped in, talking and laughing as he approached.}
}
{
}

\Verse{229}
{
    \zA{正说着，小厮来回：“冯大爷来了。”}
}
{
    \eA{The company speedily rose from table and offered him a seat. "That's right!"
        smiled Feng Tzu-ying. "You don't go out of doors,}
}
{
}

\Verse{230}
{
    \zA{宝玉便知是神武将军冯唐之子冯紫英来 了。}
}
{
    \eA{but remain at home and go in for high fun!" Both Pao-yü and Hsüeh P'an put
        on a smile. "We haven't," they remarked,}
}
{
}

\Verse{231}
{
    \zA{薛蟠等一齐都叫“快请”。}
}
{
    \eA{"seen you for ever so long. Is your venerable father strong and hale?"}
}
{
}

\Verse{232}
{
    \zA{说犹未了，只见冯紫英一路说笑已进来了，众人忙 起席让坐。}
}
{
    \eA{"My father," rejoined Tzu-ying, "is, thanks to you, strong and hale; but my
        mother recently contracted a sudden chill and has been unwell for a couple
        of days."}
}
{
}

\Verse{233}
{
    \zA{冯紫英笑道：“好啊!也不出门了，在家里高乐罢。”}
}
{
    \eA{Hsüeh P'an discerned on his face a slight bluish wound. "With whom have you
        again been boxing," he laughingly inquired, "that you've hung up this sign
        board?"}
}
{
}

\Verse{234}
{
    \zA{宝玉薛蟠都笑道： “一向少会。}
}
{
    \eA{"Since the occasion," laughed Feng Tzu-ying, "on which I wounded
        lieutenant-colonel Ch'ou's son,}
}
{
}

\Verse{235}
{
    \zA{老世伯身上安好？”}
}
{
    \eA{I've borne the lesson in mind, and never lost my temper.}
}
{
}

\Verse{236}
{
    \zA{紫英答道：“家父倒也托庇康健。}
}
{
    \eA{So how is it you say that I've again been boxing? This thing on my face was
        caused,}
}
{
}

\Verse{237}
{
    \zA{但近来家母偶 着了些风寒，不好了两天。”}
}
{
    \eA{when I was out shooting the other day on the T'ieh Wang hills, by a flap
        from the wing of the falcon."}
}
{
}

\Verse{238}
{
    \zA{薛蟠见他面上有些青伤，便笑道：“这脸上又和谁挥 拳来，挂了幌子了？”}
}
{
    \eA{"When was that?" asked Pao-yü. "I started," explained Tzu-ying, "on the 28th
        of the third moon and came back only the day before yesterday." "It isn't to
        be wondered at then,"}
}
{
}

\Verse{239}
{
    \zA{冯紫英笑道：“从那一遭把仇都尉的儿子打伤了，我记了， 再不怄气，如何又挥拳？}
}
{
    \eA{observed Pao-yü, "that when I went the other day, on the third and fourth,
        to a banquet at friend Shen's house, I didn't see you there. Yet I meant to
        have inquired about you; but I don't know how it slipped from my memory.}
}
{
}

\Verse{240}
{
    \zA{这脸上是前日打围，在铁网山叫兔鹘梢了一翅膀。”}
}
{
    \eA{Did you go alone, or did your venerable father accompany you?" "Of course,
        my father went," Tzu-ying replied, "so I had no help but to go.}
}
{
}

\Verse{241}
{
    \zA{宝玉 道：“几时的话？”}
}
{
    \eA{For is it likely, forsooth, that I've gone mad from lack of anything to do!}
}
{
}

\Verse{242}
{
    \zA{紫英道：“三月二十八日去的，前儿也就回来了。”}
}
{
    \eA{Don't we, a goodly number as we are, derive enough pleasure from our
        wine-bouts and plays that I should go in quest of such kind of fatiguing
        recreation!}
}
{
}

\Verse{243}
{
    \zA{宝玉道： “怪道前儿初三四儿我在沈世兄家赴席不见你呢!我要问，不知怎么忘了。}
}
{
    \eA{But in this instance a great piece of good fortune turned up in evil
        fortune!" Hsüeh P'an and his companions noticed that he had finished his
        tea. "Come along," they one and all proposed,}
}
{
}

\Verse{244}
{
    \zA{单你去 了，还是老世伯也去了？”}
}
{
    \eA{"and join the banquet; you can then quietly recount to us all your
        experiences."}
}
{
}

\Verse{245}
{
    \zA{紫英道：“可不是家父去!我没法儿，去罢了。}
}
{
    \eA{At this suggestion Feng Tzu-ying there and then rose to his feet. "According
        to etiquette,"}
}
{
}

\Verse{246}
{
    \zA{难道我 闲疯了，咱们几个人吃酒听唱的不乐，寻那个苦恼去？}
}
{
    \eA{he said. "I should join you in drinking a few cups; but to-day I have still
        a very urgent matter to see my father about on my return so that I truly
        cannot accept your invitation."}
}
{
}

\Verse{247}
{
    \zA{这一次，大不幸之中却有大 幸。”}
}
{
    \eA{Hsüeh P'an, Pao-yü and the other young fellows would on no account listen to
        his excuses.}
}
{
}

\Verse{248}
{
    \zA{薛蟠众人见他吃完了茶，都说道：“且入席，有话慢慢的说。”}
}
{
    \eA{They pulled him vigorously about and would not let him go. "This is, indeed,
        strange!" laughed Feng Tzu-ying. "When have you and I had,}
}
{
}

\Verse{249}
{
    \zA{冯紫英听说， 便立起身来说道：“论理，我该陪饮几杯才是，只是今儿有一件很要紧的事，回去 还要见家父面回，实不敢领。”}
}
{
    \eA{during all these years, to have recourse to such proceedings! I really am
        unable to comply with your wishes. But if you do insist upon making me have
        a drink, well, then bring a large cup and I'll take two cups full and
        finish." After this rejoinder, the party could not but give in.}
}
{
}

\Verse{250}
{
    \zA{薛蟠宝玉众人那里肯依，死拉着不放。}
}
{
    \eA{Hsüeh P'an took hold of the kettle, while Pao-yü grasped the cup, and they
        poured two large cups full.}
}
{
}

\Verse{251}
{
    \zA{冯紫英笑道： “这又奇了。}
}
{
    \eA{Feng Tzu-ying stood up and quaffed them with one draught. "But do, after
        all," urged Pao-yü,}
}
{
}

\Verse{252}
{
    \zA{你我这些年，那一回有这个道理的？}
}
{
    \eA{"finish this thing about a piece of good fortune in the midst of misfortune
        before you go." "To tell you this to-day,"}
}
{
}

\Verse{253}
{
    \zA{实在不能遵命。}
}
{
    \eA{smiled Feng Tzu-ying, "will be no great fun.}
}
{
}

\Verse{254}
{
    \zA{若必定叫我喝， 拿大杯来，我领两杯就是了。”}
}
{
    \eA{But for this purpose I intend standing a special entertainment, and inviting
        you all to come and have a long chat;}
}
{
}

\Verse{255}
{
    \zA{众人听说，只得罢了，薛蟠执壶，宝玉把盏，斟了 两大海。}
}
{
    \eA{and, in the second place, I've also got a favour to ask of you." Saying
        this, he pushed his way and was going off at once, when Hsüeh P'an
        interposed.}
}
{
}

\Verse{256}
{
    \zA{那冯紫英站着，一气而尽。}
}
{
    \eA{"What you've said," he observed, "has put us more than ever on pins and
        needles.}
}
{
}

\Verse{257}
{
    \zA{宝玉道：“你到底把这个‘不幸之幸’说完了 再走。”}
}
{
    \eA{We cannot brook any delay. Who knows when you will ask us round; so better
        tell us, and thus avoid keeping people in suspense!"}
}
{
}

\Verse{258}
{
    \zA{冯紫英笑道：“今儿说的也不尽兴，我为这个，还要特治一个东儿，请你 们去细谈一谈；二则还有奉恳之处。”}
}
{
    \eA{"The latest," rejoined Feng Tzu-ying, "in ten days; the earliest in eight."
        With this answer he went out of the door, mounted his horse, and took his
        departure. The party resumed their seats at table.}
}
{
}

\Verse{259}
{
    \zA{说着撒手就走。}
}
{
    \eA{They had another bout, and then eventually dispersed.}
}
{
}

\Verse{260}
{
    \zA{薛蟠道：“越发说的人热剌 剌的扔不下，多早晚才请我们？}
}
{
    \eA{Pao-yü returned into the garden in time to find Hsi Jen thinking with
        solicitude that he had gone to see Chia Cheng and wondering whether it
        foreboded good or evil. As soon as she perceived Pao-yü come back in a
        drunken state, she felt urged to inquire the reason of it all.}
}
{
}

\Verse{261}
{
    \zA{告诉了也省了人打闷雷。”}
}
{
    \eA{Pao-yü told her one by one the particulars of what happened. "People," added
        Hsi Jen,}
}
{
}

\Verse{262}
{
    \zA{冯紫英道：“多则十日， 少则八天。”}
}
{
    \eA{"wait for you with lacerated heart and anxious mind, and there you go and
        make merry; yet you could very well,}
}
{
}

\Verse{263}
{
    \zA{一面说，一面出门上马去了。}
}
{
    \eA{after all, have sent some one with a message." "Didn't I purpose sending a
        message?"}
}
{
}

\Verse{264}
{
    \zA{众人回来，依席又饮了一回方散。}
}
{
    \eA{exclaimed Pao-yü. "Of course, I did! But I failed to do so, as on the
        arrival of friend Feng,}
}
{
}

\Verse{265}
{
    \zA{宝玉回至园中，袭人正惦记他去见贾政，不知是祸是福，只见宝玉醉醺醺回来， 因问其原故，宝玉一一向他说了。}
}
{
    \eA{I got so mixed up that the intention vanished entirely from my mind." While
        excusing himself, he saw Pao-ch'ai enter the apartment. "Have you tasted any
        of our new things?" she asked, a smile curling her lips. "Cousin," laughed
        Pao-yü,}
}
{
}

\Verse{266}
{
    \zA{袭人道：“人家牵肠挂肚的等着，你且高乐去， 也到底打发个人来给个信儿！”}
}
{
    \eA{"you must have certainly tasted what you've got in your house long before
        us." Pao-ch'ai shook her head and smiled. "Yesterday," she said, "my brother
        did actually make it a point to ask me to have some;}
}
{
}

\Verse{267}
{
    \zA{宝玉道：“我何尝不要送信儿，因冯世兄来了，就 混忘了。”}
}
{
    \eA{but I had none; I told him to keep them and send them to others, so
        confident am I that with my mean lot and scanty blessings I little deserve
        to touch such dainties."}
}
{
}

\Verse{268}
{
    \zA{正说着，只见宝钗走进来，笑道：“偏了我们新鲜东西了。”}
}
{
    \eA{As she spoke, a servant-girl poured her a cup of tea and brought it to her.
        While she sipped it, she carried on a conversation on irrelevant matters;}
}
{
}

\Verse{269}
{
    \zA{宝玉笑道： “姐姐家的东西，自然先偏了我们了。”}
}
{
    \eA{which we need not notice, but turn our attention to Lin Tai-yü. The instant
        she heard that Chia Cheng had sent for Pao-yü,}
}
{
}

\Verse{270}
{
    \zA{宝钗摇头笑道：“昨儿哥哥倒特特的请我 吃，我不吃，我叫他留着送给别人罢。}
}
{
    \eA{and that he had not come back during the whole day, she felt very distressed
        on his account. After supper, the news of Pao-yü's return reached her, and
        she keenly longed to see him and ask him what was up.}
}
{
}

\Verse{271}
{
    \zA{我知道我的命小福薄，不配吃那个。”}
}
{
    \eA{Step by step she trudged along, when espying Pao-ch'ai going into Pao-yü's
        garden, she herself followed close in her track.}
}
{
}

\Verse{272}
{
    \zA{说着， 丫鬟倒了茶来，吃茶说闲话儿，不在话下。}
}
{
    \eA{But on their arrival at the Hsin Fang bridge, she caught sight of the
        various kinds of water-fowl, bathing together in the pond,}
}
{
}

\Verse{273}
{
    \zA{却说那黛玉听见贾政叫了宝玉去了，一日不回来，心中也替他忧虑。}
}
{
    \eA{and although unable to discriminate the numerous species, her gaze became so
        transfixed by their respective variegated and bright plumage and by their
        exceptional beauty,}
}
{
}

\Verse{274}
{
    \zA{至晚饭后， 闻得宝玉来了，心里要找他问问是怎么样了，一步步行来。}
}
{
    \eA{that she halted. And it was after she had spent some considerable time in
        admiring them that she repaired at last to the I Hung court. The gate was
        already closed. Tai-yü, however, lost no time in knocking. But Ch'ing Wen
        and Pi Hen had,}
}
{
}

\Verse{275}
{
    \zA{见宝钗进宝玉的园内去 了，自己也随后走了来。}
}
{
    \eA{who would have thought it, been having a tiff, and were in a captious mood,
        so upon unawares seeing Pao-ch'ai step on the scene,}
}
{
}

\Verse{276}
{
    \zA{刚到了沁芳桥，只见各色水禽尽都在池中浴水，也认不出 名色来，但见一个个文彩？}
}
{
    \eA{Ch'ing Wen at once visited her resentment upon Pao-ch'ai. She was just
        standing in the court giving vent to her wrongs, shouting: "You're always
        running over and seating yourself here, whether you've got good reason for
        doing so or not;}
}
{
}

\Verse{277}
{
    \zA{灼，好看异常，因而站住，看了一回。}
}
{
    \eA{and there's no sleep for us at the third watch, the middle of the night
        though it be,"}
}
{
}

\Verse{278}
{
    \zA{再往怡红院来， 门已关了，黛玉即便叩门。}
}
{
    \eA{when, all of a sudden, she heard some one else calling at the door. Ch'ing
        Wen was the more moved to anger. Without even asking who it was,}
}
{
}

\Verse{279}
{
    \zA{谁知晴雯和碧痕二人正拌了嘴，没好气，忽见宝钗来了， 那晴雯正把气移在宝钗身上，偷着在院内抱怨说：“有事没事跑了来坐着，叫我们 三更半夜的不得睡觉！”}
}
{
    \eA{she rapidly bawled out: "They've all gone to sleep; you'd better come
        to-morrow." Lin Tai-yü was well aware of the natural peculiarities of the
        waiting-maids, and of their habit of playing practical jokes upon each
        other, so fearing that the girl in the inner room had failed to recognise
        her voice, and had refused to open under the misconception that it was some
        other servant-girl,}
}
{
}

\Verse{280}
{
    \zA{忽听又有人叫门，晴雯越发动了气，也并不问是谁，便说 道：“都睡下了，明儿再来罢！”}
}
{
    \eA{she gave a second shout in a higher pitch. "It's I!" she cried, "don't you
        yet open the gate?" Ch'ing Wen, as it happened, did not still distinguish
        her voice; and in an irritable strain,}
}
{
}

\Verse{281}
{
    \zA{黛玉素知丫头们的性情，他们彼此玩耍惯了，恐怕院内的丫头没听见是他的声 音，只当别的丫头们了，所以不开门；因而又高声说道：“是我，还不开门么？”}
}
{
    \eA{she rejoined: "It's no matter who you may be; Mr. Secundus has given orders
        that no one at all should be allowed to come in." As these words reached Lin
        Tai-yü's ear, she unwittingly was overcome with indignation at being left
        standing outside. But when on the point of raising her voice to ask her one
        or two things,}
}
{
}

\Verse{282}
{
    \zA{晴雯偏偏还没听见，便使性子说道：“凭你是谁，二爷吩咐的，一概不许放进人来 呢！”}
}
{
    \eA{and to start a quarrel with her; "albeit," she again argued mentally, "I can
        call this my aunt's house, and it should be just as if it were my own, it's,
        after all, a strange place,}
}
{
}

\Verse{283}
{
    \zA{黛玉听了这话，不觉气怔在门外。}
}
{
    \eA{and now that my father and mother are both dead, and that I am left with no
        one to rely upon,}
}
{
}

\Verse{284}
{
    \zA{待要高声问他，逗起气来，自己又回思一 番：“虽说是舅母家如同自己家一样，到底是客边。}
}
{
    \eA{I have for the present to depend upon her family for a home. Were I now
        therefore to give way to a regular fit of anger with her, I'll really get no
        good out of it." While indulging in reflection, tears trickled from her
        eyes.}
}
{
}

\Verse{285}
{
    \zA{如今父母双亡，无依无靠，现 在他家依栖，若是认真怄气，也觉没趣。”}
}
{
    \eA{But just as she was feeling unable to retrace her steps, and unable to
        remain standing any longer, and quite at a loss what to do, she overheard
        the sound of jocular language inside,}
}
{
}

\Verse{286}
{
    \zA{一面想，一面又滚下泪珠来了。}
}
{
    \eA{and listening carefully, she discovered that it was, indeed, Pao-yü and
        Pao-ch'ai.}
}
{
}

\Verse{287}
{
    \zA{真是回 去不是，站着不是。}
}
{
    \eA{Lin Tai-yü waxed more wroth. After much thought and cogitation,}
}
{
}

\Verse{288}
{
    \zA{正没主意，只听里面一阵笑语之声，细听一听，竟是宝玉宝钗 二人。}
}
{
    \eA{the incidents of the morning flashed unawares through her memory. "It must,
        in fact," she mused, "be because Pao-yü is angry with me for having
        explained to him the true reasons.}
}
{
}

\Verse{289}
{
    \zA{黛玉心中越发动了气，左思右想，忽然想起早起的事来：“必竟是宝玉恼我 告他的原故。}
}
{
    \eA{But why did I ever go and tell you? You should, however, have made inquiries
        before you lost your temper to such an extent with me as to refuse to let me
        in to-day; but is it likely that we shall not by and bye meet face to face
        again?"}
}
{
}

\Verse{290}
{
    \zA{但只我何尝告你去了？}
}
{
    \eA{The more she gave way to thought, the more she felt wounded and agitated;}
}
{
}

\Verse{291}
{
    \zA{你也不打听打听，就恼我到这步田地!你今儿不 叫我进来，难道明儿就不见面了？”}
}
{
    \eA{and without heeding the moss, laden with cold dew, the path covered with
        vegetation, and the chilly blasts of wind, she lingered all alone, under the
        shadow of the bushes at the corner of the wall,}
}
{
}

\Verse{292}
{
    \zA{越想越觉伤感，便也不顾苍苔露冷，花径风寒， 独立墙角边花阴之下，悲悲切切，呜咽起来。}
}
{
    \eA{so thoroughly sad and dejected that she broke forth into sobs. Lin Tai-yü
        was, indeed, endowed with exceptional beauty and with charms rarely met with
        in the world. As soon therefore as she suddenly melted into tears,}
}
{
}

\Verse{293}
{
    \zA{原来这黛玉秉绝代之姿容，具稀世之 俊美，不期这一哭，把那些附近的柳枝花朵上宿鸟栖鸦，一闻此声，俱忒楞楞飞起 远避，不忍再听。}
}
{
    \eA{and the birds and rooks roosting on the neighbouring willow boughs and
        branches of shrubs caught the sound of her plaintive tones, they one and all
        fell into a most terrific flutter, and, taking to their wings, they flew
        away to distant recesses, so little were they able to listen with equanimity
        to such accents.}
}
{
}

\Verse{294}
{
    \zA{正是。}
}
{
    \eA{But the spirits of the flowers were,}
}
{
}

\Verse{295}
{
    \zA{花魂点点无情绪，鸟梦痴痴何处惊。}
}
{
    \eA{at the time, silent and devoid of feeling, the birds were plunged in dreams
        and in a state of stupor,}
}
{
}

\Verse{296}
{
    \zA{因又有一首诗道： 颦儿才貌世应稀，独抱幽芳出绣闺。}
}
{
    \eA{so why did they start? A stanza appositely assigns the reason:— P'in Erh's
        mental talents and looks must in the world be rare—. Alone, clasped in a
        subtle smell,}
}
{
}

\Verse{297}
{
    \zA{呜咽一声犹未了，落花满地鸟惊飞。}
}
{
    \eA{she quits her maiden room. The sound of but one single sob scarcely dies
        away,}
}
{
}

\Verse{298}
{
    \zA{那黛玉正自啼哭，忽听吱娄娄一声，院门开处，不知是那一个出来。}
}
{
    \eA{And drooping flowers cover the ground and birds fly in dismay. Lin Tai-yü
        was sobbing in her solitude, when a creaking noise struck her ear and the
        door of the court was flung open.}
}
{
}

\Verse{299}
{
    \zA{要知端的，下回分解。}
}
{
    \eA{Who came out, is not yet ascertained; but, reader, should you wish to know,}
}
{
}

\Verse{300}
{
    \zA{(本章完)}
}
{
    \eA{the next chapter will explain.}
}
{
}
